\subsection{Derived exponential comparisons and the remaining numerical lift}

\paragraph{D-024. Known, with checked scope: Zakharov's comparison.}
Zakharov's Theorem~3.3.6, printed p.~10, compares derived exponential
functors on the duals of perfect complexes in $D^{<-1}(\mathsf{Ab})$;
the target is the ordinary derived category. Applying it to binomial
and divided-power functors gives the comparison in Section~6,
pp.~17--18. Proposition~5.0.3, p.~15, distinguishes the functorial
identity $\operatorname{gr}\operatorname{Bin}\cong\Gamma$ from the
coalgebra identification $\operatorname{Bin}''\cong\Gamma''$ on the
restricted category $\mathsf{Ab}''$. The decreasing numerical-tail
filtration immediately following that proposition is explicitly not
functorial for arbitrary linear maps. Remark~4.0.9, p.~14, also
explicitly excludes an extension of the stated Dold--Puppe--Thom
splitting to the infinity-categorical localization.
These are restrictions in the source, rather than an assertion that
all conceivable enhanced comparisons fail \cite{D_Zakharov2026}.

The original two-term input $M_n=[\Phi_n\to\Psi_n]$ has cohomology
$\Z T$ in degree zero, and therefore does not satisfy that
connectivity hypothesis. Applying the derived binomial functor
without shifting does change the exterior rows: its weight-two
cosimplicial normalization has the additional $\Gamma^2\Psi$ and
$\Psi\otimes\Psi$ terms displayed in Statement~\ref{sa08:A022}.
Consequently it evades the hypotheses of Statement~\ref{sa08:A024}
by changing those rows, but its underived sections already give the
wrong cohomology on the finite test space in Statement~\ref{sa08:A022}.
Nor can a uniform positive shift followed by the evident weight
deshearing repair the associated graded, as follows.

\begin{proposition}[D-025: \status{Proved}; every positive uniform shift fails]
\label{sa08:D025}
Use cohomological shifts, so $\Z[-r]$ is concentrated in degree $r$.
For $r\ge1$, put
\[
 E_r^d=\mathbf L\Gamma^d(\Z[-r])[rd].
\]
The desheared graded object $\bigoplus_{d\ge0}E_r^d$ cannot be the
graded divided-power algebra $\Gamma_{\Z}(\Z T)$ concentrated in
cohomological degree zero. For $r=1$ its weight-two piece is zero.
For $r=2$ it has the polynomial-algebra multiplication rather than
the divided-power multiplication. For every $r\ge3$ its weight-two
piece has nonzero $2$-torsion in cohomology.
Thus no such shift repairs the numerical construction even at the
single vertex $M_0=\Z T$.
\end{proposition}

\begin{proof}
For $r=1$, d\'ecalage gives
$\mathbf L\Gamma^2(\Z[-1])\simeq\Lambda^2\Z[-2]=0$.
For $r=2$, multiplicative double d\'ecalage gives
\[
 \bigoplus_d H^{2d}\bigl(\mathbf L\Gamma^d(\Z[-2])\bigr)
   \cong\operatorname{Sym}_{\Z}(\Z).
\]
The weight-one generator therefore has square a generator of
weight two. In $\Gamma_{\Z}(\Z T)$ its square is $2T^{[2]}$.
These graded algebras cannot be isomorphic, as is also seen after
reduction modulo two. The multiplicative d\'ecalage used here is
stated in Zakharov, Theorem~7.1.1, p.~18.

Here is the complete additive calculation for larger shifts. Set
$C_r=\mathbf L\Gamma^2(\Z[-r])$. Perfect duality identifies it with
\[
 C_r\simeq
 \bigl(\mathbf L\operatorname{Sym}^2(\Z[r])\bigr)^\vee.
 \tag{D-25a}
\]
For example this duality is recorded in
\cite[Remark~2.70]{A_KSZ}, version~2.
Write $G_n$ for the homological complex
$N_s\Gamma^2(K(\Z,n))$. Double d\'ecalage says that
$\mathbf L\operatorname{Sym}^2(\Z[r])$ has homology
$H_{j-4}(G_{r-2})$ in homological degree $j$, for $r\ge2$.
The precise primary references are
\cite[Proposition~3.5, equation~(3-9), p.~270; Proposition~8.7,
equation~(8-16), p.~306]{D_BMT2016}.

For clarity, the rank-one calculation can also be obtained
recursively. On free abelian groups there is an exact sequence
\[
 0\longrightarrow\operatorname{Sym}^2(V)
 \longrightarrow\Gamma^2(V)
 \longrightarrow V/2V\longrightarrow0,
\]
where the first map sends $v^2$ to $2\gamma_2(v)$ and the last
map sends $\gamma_2(v)$ to $\bar v$. Applied to $K(\Z,n)$,
this sequence and double d\'ecalage give, for $n\ge2$, a new
$\Z/2$ in degree $n$, with the remaining homology equal to that
of $G_{n-2}[4]$. The initial values are $G_0\simeq\Z$ and
$H_1(G_1)=\Z/2$, with no other homology for $G_1$.
The degree ranges are disjoint, so the long exact sequence has no
extension ambiguity. It follows that $G_n$ has a $\Z/2$ in
degrees $n,n+2,\ldots,2n-2$ for even $n$, and in degrees
$n,n+2,\ldots,2n-1$ for odd $n$; in addition it has $\Z$ in
degree $2n$ when $n$ is even.

The complexes in (D-25a) are perfect. Integral duality sends a
free homology class in degree $j$ to cohomological degree $j$,
and a $\Z/2$ class to degree $j+1$. Thus for even $r\ge2$,
\[
 H^j(C_r)=
 \begin{cases}
 \Z,&j=2r,\\
 \Z/2,&j=r+3,r+5,\ldots,2r-1,\\
 0,&\text{otherwise};
 \end{cases}
\]
the middle range is empty for $r=2$. For odd $r\ge3$,
\[
 H^j(C_r)=
 \begin{cases}
 \Z/2,&j=r+3,r+5,\ldots,2r,\\
 0,&\text{otherwise}.
 \end{cases}
\]
Shifting by $[2r]$ leaves these nonzero torsion groups present.
They cannot agree with the free rank-one weight-two piece of
$\Gamma_{\Z}(\Z T)$.
\end{proof}

The saved certificate \texttt{code/D\_shift\_quadratic\_checks.py}
constructs the actual integral normalized face matrices of
$\operatorname{Sym}^2K(\Z,r)$ and $\Gamma^2K(\Z,r)$ for
$1\le r\le4$. Its 50 exact assertions check the differential
identities, Smith-form homology, and the unit Alexander--Whitney
coefficient giving the shift-two square. These finite checks
supplement the all-index proof above.

\begin{proposition}[D-026: \status{Proved}; a total derived lift and its limitation]
\label{sa08:D026}
Let $R=\operatorname{Num}(\Z T)$ with its increasing degree
filtration. The constant simplicial filtered $E_\infty$ algebra
$R$ has associated graded equivalent to the PD de Rham algebra
$A_\bullet$ in
$\operatorname{Fun}(\Delta^{\mathrm{op}},
\operatorname{CAlg}(D(\Z)))$, with its weight grading retained.
For a simplicial set $X$, define the filtered algebra $B_X$ by
\[
 F_wB_X=R\Gamma(X,R_{\le w}),\qquad
 B_X=\operatorname*{colim}_{w\ge0}F_wB_X.
 \tag{D-26a}
\]
Then $\operatorname{gr}_wB_X\simeq R\Gamma(X,\Z)$ for every
$w\ge0$, and its associated graded recovers the PD cochains
after the shadow quotient:
\[
 Q\operatorname{gr}B_X
   \simeq Q C_{\mathrm{PD},\mathrm{gr}}(X)
\]
as enhanced commutative algebra objects. Before this quotient,
the associated graded is not the weighted underived PD cochain
object in general.
\end{proposition}

\begin{proof}
The associated graded of the degree filtration on $R$ is $S$.
The canonical inclusion of constant coefficients
$S\to A_n$ is an objectwise Poincar\'e quasi-isomorphism,
compatible with products, weight, and all simplicial operators.
It is therefore an equivalence of the indicated diagrams of
$E_\infty$ algebras.

Derived sections are lax symmetric monoidal. The coefficient
multiplications $R_{\le a}\otimes R_{\le b}\to R_{\le a+b}$
give (D-26a) its filtered $E_\infty$ algebra structure.
The exact coefficient sequence with quotient
$R_{\le w}/R_{\le w-1}\cong\Z$ remains a cofiber sequence
after derived sections, because limits in a stable category
preserve finite cofiber sequences. This proves the assertion
about each associated-graded piece, including its
divided-power coefficient multiplication. The enhanced shadow
comparison is precisely Statement D-006E.

For $X=S^1$, weight zero on the left has
$H^1=\Z$, whereas the literal weight-zero PD cochains have
no positive cohomology. Hence this is a shadow comparison,
not an all-weight comparison before $Q$.
For infinite $X$, the colimit in (D-26a) must be retained:
the degree filtration need not exhaust all $R$-valued
cochains, so one must not silently replace this colimit by
$R\Gamma(X,R)$. For finite simplicial $X$ the usual
identification does hold.
\end{proof}

\paragraph{The minimal additional obstruction data.}
Statement~\ref{sa08:D026} prevents a blanket nonexistence theorem
for arbitrary total derived lifts. A model becomes subject to
Statement~\ref{sa08:A024} once it also retains a relative
exterior-degree-zero realization $\operatorname{Num}(\Phi)$,
a coherent module first row with horizontal homology $\Z[1]$
and the augmentation action of $R$, and a first differential
with nonzero $T$-transgression. Already the corresponding
morphism $\Z\to\Z[1]$ in $D(R)$ is impossible. No higher
operadic structure can change this zero Ext group while
preserving those data. An actual all-$X$ filtered lift of the
weighted cochains which does not retain this realization is a
distinct question; the constant construction does not answer it.
